\subsection{The Leray--Schauder theorem}

The following theorem will be proved in TA, forthcoming.

THEOREM 1 (Brouwer). --- Let B be a nonempty compact convex subset of a finite-dimensional normed vector space. Every continuous map from B into B has a fixed point.

We shall deduce from this the following result.

THEOREM 2 (Leray--Schauder). --- Let X be a nonempty compact topological space, homeomorphic to a convex subset of a locally convex space. Every continuous map from X into X has a fixed point.

We may assume that X is a convex subset of a locally convex space E. Let N be the closure of \(\{0\}\) in E and \(\pi: E \rightarrow E/N\) the canonical surjection. Since X is a separated space, the restriction of \(\pi\) to X is injective; since X is compact, it then defines a homeomorphism from X onto a nonempty compact convex subset of the separated locally convex space E/N. This allows us to assume that the space E is separated.

Let \(f: \mathrm{X} \to \mathrm{X}\) be a continuous map. The map \(h: \mathrm{X} \to \mathrm{E}\) defined by \(h(x) = f(x) - x\) is continuous; since X is compact, the image \(h(\mathrm{X})\) is closed. It therefore suffices to prove that 0 is adherent to \(h(\mathrm{X})\) , that is, for every open convex neighborhood U of 0 in E, there exists a point \(b\) of X such that \(f(b) - b \in \mathrm{U}\) .

Let U be a convex neighborhood of 0 in E. For every \(a \in X\) , we denote by \(V_a\) the set of elements \(x\) of X such that \(f(x) - f(a) \in U\) . It is open in X and contains \(a\) . Since X is compact and nonempty, there exists a nonempty finite subset A of X such that the sets \(V_a\) , for \(a \in A\) , cover X. Let \((\varphi_a)_{a \in A}\) be a continuous partition of unity subordinate to the covering \((V_a)_{a \in A}\) (TG, IX, p. 43, prop. 1 and p. 47, th. 3). Let F denote the subspace of E generated by \(f(A)\) . For every \(x \in F \cap X\) set

\[
g (x) = \sum_ {a \in A} \varphi_ {a} (x) f (a).
\]

One thus defines a continuous map \(g: F \cap X \to E\) . Its image is contained in the convex hull of \(f(A)\) , hence in \(F \cap X\) . Since the set \(F \cap X\) is a nonempty compact convex subset of a finite-dimensional vector space, the map g has a fixed point \(b \in F \cap X\) by Theorem 1. We have

\[
f (b) - b = f (b) - g (b) = \sum_ {a \in \mathrm{A}} \varphi_ {a} (b) (f (b) - f (a)).
\]

For every \(a \in \mathrm{A}\) such that \(\varphi_{a}(b) \neq 0\) , we have \(b \in \mathrm{V}_{a}\) so that \(f(b) - f(a) \in \mathrm{U}\) . Since U is convex, we conclude that \(f(b) - b \in \mathrm{U}\) , which completes the proof.

